\subsection{The actual arithmetic normal contribution is strictly positive}

We now use the actual zero set of the amplitude in (FC.1).
Hardy's theorem that infinitely many zeros of the Riemann zeta
function lie on the critical line, with the original source and
reading record retained in Section~\ref{sec:theta-gamma-reference},
implies that
\[
 \mathcal Z_h=\{T\in\mathbb R:a_h^{\rm SC}(T)=0\}
 \quad\hbox{is infinite and discrete.}
 \tag{FC.21}
\]
Indeed a finite polynomial \(h\) removes only its selected finite
set of complete zero orders. Every remaining critical-line zero is
still a zero of \(v_h\), and the nonzero entire function \(v_h\)
has a discrete zero set. This also applies to \(h=1\).
At any \(T\in\mathcal Z_h\), keep its entire order
\(\ell_T\ge1\) and its original leading coefficient
\[
 a_h^{\rm SC}(T+\delta)
       =c_T\delta^{\ell_T}+O(\delta^{\ell_T+1}),\qquad
 (a_h^{\rm SC})'(T+\delta)
       =\ell_Tc_T\delta^{\ell_T-1}+O(\delta^{\ell_T}),
 \quad c_T=\frac{(a_h^{\rm SC})^{(\ell_T)}(T)}{\ell_T!}\ne0.
 \tag{FC.22}
\]
These are expansions of the original amplitude, including its fixed
phase and its factor \(1/\sqrt{2\pi}\).

\begin{theorem}[Strict scalar residual and every finite normal direction]
For the actual amplitude and every integer \(k\ge2\),
\[
 \int_{\mathbb R}\|n(u)\|^2du>0,\qquad
 \boxed{\int_{\mathbb R}\frac{|m'(u)|^2}{m(u)}du
                <\frac{\mu_h^{k-1}}k I_h.}
 \tag{FC.23}
\]
For \(k\ge3\), \(\|n(u)\|>0\) for every real \(u\).
For any nonempty finite independent relative polynomial frame of
(FC.5), its full normal Gram satisfies
\[
 \boxed{\begin{aligned}
  \mathcal N(u)&>0&&\text{for every }u\in\mathbb R,
                                      &&k\ge3,\\
  \mathcal N(u)&>0&&\text{for every }u\notin
                  \mathcal Z_h+\mathcal Z_h,
                                      &&k=2.
 \end{aligned}}
 \tag{FC.24}
\]
In particular \(\mathcal N(u)>0\) almost everywhere for every
\(k\ge2\), since \(\mathcal Z_h+\mathcal Z_h\) is countable.
No lower bound uniform in \(u\), the relative degree, or the
packet is asserted by these strict inequalities.
\end{theorem}

\begin{proof}
First fix \(T\in\mathcal Z_h\) and choose
\(b_2,\ldots,b_k\) outside \(\mathcal Z_h\). Put
\(u_0=T+\sum_{i=2}^kb_i\) and
\(A_b=\prod_{i=2}^ka_h^{\rm SC}(b_i)\ne0\).
Within that fixed original sum fibre consider the transverse path
\[
 t_1=T+\delta,\quad t_i=b_i\ (2\le i<k),\quad
 t_k=b_k-\delta.
\]
It changes the relative variables while retaining \(u=u_0\).
The derivative \(\partial_u\) is still the original derivative
at fixed relative variables in (FC.2), and its product formula gives
\[
 \begin{aligned}
  \Psi(u_0,\mathbf y(\delta))
      &=c_TA_b\delta^{\ell_T}+O(\delta^{\ell_T+1}),\\
  \partial_u\Psi(u_0,\mathbf y(\delta))
      &=\frac{\ell_T}{k}c_TA_b\delta^{\ell_T-1}
                              +O(\delta^{\ell_T}).
 \end{aligned}
 \tag{FC.25}
\]
Only the derivative of the first amplitude has order
\(\ell_T-1\); every other product term retains the full
factor \(a_h^{\rm SC}(T+\delta)\) and has order at least
\(\ell_T\). Since \(m(u_0)>0\) and \(m\) is smooth,
subtracting \(m'(u_0)\Psi/(2m(u_0))\) does not change that
nonzero leading coefficient. Thus \(n(u_0,\cdot)\) is nonzero
at nearby points of this path and, by continuity, on a nonempty
open subset of the entire relative fibre. Its squared fibre norm
is strictly positive. Differentiated Schwartz estimates show that
\(\Psi\) and \(\partial_u\Psi\) are continuous as functions
of \(u\) with values in \(\mathscr H\). Smoothness and positivity
of \(m\) therefore make \(n\) continuous in that Hilbert norm.
Its squared norm stays positive on some interval around \(u_0\),
proving positive integrated residual. Equation (FC.4) proves (FC.23).

For \(k\ge3\), the same construction works at every prescribed
\(u_0\). Choose \(b_3,\ldots,b_{k-1}\) outside the zero set
if that list is nonempty. Choose \(b_2\) outside both
\(\mathcal Z_h\) and the translate-reflection of that set which
would make
\(b_k=u_0-T-\sum_{i=2}^{k-1}b_i\) a zero.
These two discrete sets cannot exhaust the real line. This proves
the pointwise scalar statement.

To prove the full matrix assertion, fix \(u\) and suppose
\(N_uc=0\). The exact projection formula gives
\(j_u'c=j_ud\), where \(d=\Gamma_uc\) is a finite coefficient
vector. Write \(\vartheta_c(\mathbf y)=
\sum_\alpha\theta_\alpha(i\mathbf y)c_\alpha\), and similarly
\(\vartheta_d\). Equality in \(\mathscr H\) says
\[
 (\partial_u\Psi)\vartheta_c=\Psi\vartheta_d.
\]
Both sides have smooth representatives in the relative variables,
so their almost-everywhere equality is equality everywhere: a nonzero
continuous difference at a point would remain nonzero on a set of
positive measure.

Suppose \(k\ge3\). For each \(T\in\mathcal Z_h\), restrict
to the affine hyperplane \(t_1=u/k+y_1=T\), of dimension
\(k-2\ge1\). Each of the other original coordinates is a
nonconstant affine function on this hyperplane. The points where
one of their amplitude factors vanishes are consequently a countable
union of proper affine hyperplanes of measure zero. Their complement
is dense and is open near each of its points. At any such point
\(p\), apply (FC.25). In the asserted smooth identity the left
side has leading coefficient
\(\ell_Tc_TA_b\vartheta_c(p)/k\) in degree \(\ell_T-1\),
whereas the right side has order at least \(\ell_T\).
Hence \(\vartheta_c(p)=0\).
Continuity proves its vanishing on the entire real hyperplane.
Division in \(y_1\) then shows that the polynomial
\(\vartheta_c\) has factor \(y_1-(T-u/k)\): the remainder
polynomial in the other variables vanishes on all their real values,
and therefore is identically zero.
There are infinitely many distinct such \(T\), while the degree
in \(y_1\) is finite. Thus \(\vartheta_c=0\), and independence
of the original relative polynomials gives \(c=0\).
This proves injectivity of \(N_u\) and therefore
\(\mathcal N(u)>0\) on the whole finite coefficient space.

For \(k=2\), take \(u\notin\mathcal Z_h+\mathcal Z_h\).
At every \(T\in\mathcal Z_h\), the other amplitude
\(a_h^{\rm SC}(u-T)\) is nonzero. The same expansion forces
the univariate polynomial \(\vartheta_c\) to vanish at
\(y_1=T-u/2\). These are infinitely many distinct points,
so again \(\vartheta_c=0\) and \(c=0\).
The discrete subset \(\mathcal Z_h\) of \(\mathbb R\)
is countable, hence so is its set of pairwise sums. This proves
every assertion in (FC.24).
\end{proof}

The consequences hold for the actual full derivative graph and both
quadratic branches. For every nonzero admitted source polynomial
\(P\),
\[
 \begin{gathered}
 0<\int_{\mathbb R}c_P(u)^*\mathcal N(u)c_P(u)\,du<\infty,\\
 \mathsf Q(x)>0\quad\text{for all }x>0\text{ when }k\ge3,
 \qquad
 \mathsf Q(x)>0\quad\text{for almost every }x>0\text{ when }k=2.
 \end{gathered}
 \tag{FC.26}
\]
Indeed a nonzero vector polynomial \(c_P\) is nonzero outside
a finite set on the real line, and (FC.24) supplies strict
positivity almost everywhere. The integral is finite because
\(\|N_uc_P\|\le\|j_u'c_P\|\), and the latter is a
polynomial multiple of differentiated Schwartz columns.
The assertion about \(\mathsf Q\) follows from its invertible
congruence (FC.14). For \(k=2\), the exceptional positive
\(x\) for which either \(\sqrt x\) or \(-\sqrt x\) belongs
to \(\mathcal Z_h+\mathcal Z_h\) form a countable set.
Under every declared smooth invertible frame,
\(N_u^Cc_P^C(u)=N_uc_P(u)\) and
\(\mathcal N^C=C^*\mathcal NC\).
Thus the full normal energy and strictness are preserved by the
exact coefficient isomorphism \(c_P^C=C^{-1}c_P\), including
its original filtered image.

The excluded two-factor fibre can actually vanish. For the quartet
packet, the original amplitude is even, so
\[
 \partial_u\Psi(0,y)=\tfrac12\bigl(
 (a_h^{\rm SC})'(y)a_h^{\rm SC}(-y)
 +a_h^{\rm SC}(y)(a_h^{\rm SC})'(-y)\bigr)=0.
 \tag{FC.27}
\]
Since the relative frame is independent of \(u\), this gives
\(j_0'=0\), \(N_0=0\), and \(\mathcal N(0)=0\).
Also \(m'(0)=0\), hence \(n(0)=0\) exactly. The positive
total residual and the almost-everywhere full matrix assertion
therefore retain this original zero fibre. They prove that every
nonzero actual finite polynomial direction has a positive normal
contribution while leaving its complete size in (FC.15) and (FC.26).
